\PassOptionsToPackage{table,xcdraw}{xcolor}
\documentclass[11pt, a4paper]{article}

\usepackage[T1]{fontenc}
\usepackage[utf8]{inputenc}
\usepackage{tgpagella}
\usepackage[spanish, es-noshorthands]{babel}
\usepackage{amsmath, amssymb}
\usepackage{booktabs}
\usepackage{tabularx}
\usepackage[most]{tcolorbox}
\usepackage[margin=2cm]{geometry}

\definecolor{ucbblue}{RGB}{0, 51, 102}

\begin{document}

\begin{center}
    \Large\textbf{\textcolor{ucbblue}{Guía Estructural: Mini-Artículo IEEE de Laboratorio}}[cite: 12]
\end{center}

\begin{tcolorbox}[colback=white, colframe=ucbblue, title=\textbf{Propósito del Informe[cite: 12]}]
    Comunicar un argumento de ingeniería coherente, transformando: Problema $\rightarrow$ Modelo $\rightarrow$ Método $\rightarrow$ Evidencia $\rightarrow$ Comparación $\rightarrow$ Diagnóstico $\rightarrow$ Conclusión[cite: 12]. \\
    \textbf{Extensión objetivo:} 3 a 4 páginas estilo IEEE por laboratorio[cite: 12].
\end{tcolorbox}

\section*{1. Estructura Requerida[cite: 12]}
\begin{enumerate}
    \setlength{\itemsep}{0pt}
    \item \textbf{Title / Abstract:} 100-150 palabras con problema, método, resultado cuantitativo y conclusión[cite: 12]. Sin teoría extensa ni resultados inventados[cite: 12].
    \item \textbf{Index Terms / Keywords}[cite: 12].
    \item \textbf{I. Introduction:} ¿Qué problema de ingeniería se investiga y por qué importa? (Fenómeno, objetivo, pregunta experimental)[cite: 12].
    \item \textbf{II. Theoretical Background:} Solo la teoría necesaria con símbolos definidos, ecuaciones numeradas, unidades y valores reales sustituidos[cite: 12].
    \item \textbf{III. Experimental Methodology:} Circuito, valores, Op-Amp exacto, fabricante, instrumentos, alimentación, variables y secuencia. La simulación va en subsección separada[cite: 12].
    \item \textbf{IV. Results:} Tablas, gráficos y métricas. Use etiquetas de evidencia: \texttt{THEORETICAL}, \texttt{SIMULATED}, \texttt{DATASHEET}, \texttt{MEASURED}[cite: 12].
    \item \textbf{V. Discussion:} Compare teoría $\leftrightarrow$ simulación $\leftrightarrow$ datasheet $\leftrightarrow$ medición[cite: 12]. Responda: ¿Coincidió? ¿Causa física plausible? ¿Explicaciones alternativas? ¿Limitaciones?[cite: 12]
    \item \textbf{VI. Conclusions:} 3-5 conclusiones concisas derivadas de la evidencia, con valor cuantitativo[cite: 12].
    \item \textbf{References:} Datasheets (prioridad 1), fuentes académicas, manuales de instrumentos[cite: 12].
\end{enumerate}

\section*{2. Política de Evidencia Faltante[cite: 12]}
Si una medición no existe, escriba explícitamente: \texttt{Measurement not obtained}[cite: 12]. \textbf{NUNCA} sustituya silenciosamente una medición por una simulación o un valor inventado[cite: 12].

\end{document}
